% บทคัดย่อภาษาไทย — ไม่เกิน 300 คำ
% ครอบคลุม: วัตถุประสงค์ → วิธีการ → ผลลัพธ์ → การประยุกต์ใช้

โครงงานนี้มีวัตถุประสงค์เพื่อเปรียบเทียบสถาปัตยกรรมการสร้างคำตอบโดยอาศัยการค้นคืนแบบมาตรฐานกับแบบที่มีเอเจนต์ตัดสินใจเอง บนระบบแนะนำสินค้าเชิงสนทนาที่รับคำถามได้ทั้งภาษาไทยและภาษาอังกฤษ โดยคลังสินค้าเป็นภาษาอังกฤษทั้งหมด

วิธีดำเนินงานเริ่มจากเตรียมคลังสินค้าจริงจาก Amazon จำนวน 81{,}264 รายการ สร้างดัชนีค้นหาเชิงความหมายด้วยแบบจำลองฝังเวกเตอร์สองภาษา แล้วพัฒนาสถาปัตยกรรม 3 แบบที่ใช้ส่วนประกอบร่วมกันทั้งหมดและต่างกันเฉพาะชั้นที่แบบจำลองภาษาตัดสินใจเอง เพื่อให้การเปรียบเทียบเปลี่ยนตัวแปรทีละตัว ระบบทั้งหมดประมวลผลบนเครื่องคอมพิวเตอร์ส่วนบุคคลด้วยแบบจำลองภาษาขนาด 8 พันล้านพารามิเตอร์ การประเมินใช้ชุดทดสอบที่สร้างขึ้นเอง 140 คำค้นสองภาษา แบ่งตามสิ่งที่คลังมี และติดเฉลยความตรงด้วยมือตามเกณฑ์มาตรฐาน 4 ระดับ

ผลการทดลองพบว่าสถาปัตยกรรมแบบเอเจนต์ให้คะแนน nDCG@3 สูงกว่าแบบมาตรฐานทั้งสองภาษาและไม่แพ้แม้แต่คำค้นเดียว โดยฝั่งอังกฤษเพิ่มจาก 0.3288 เป็น 0.3564 อย่างมีนัยสำคัญทางสถิติ และฝั่งไทยเพิ่มจาก 0.3165 เป็น 0.3267 อย่างไรก็ตาม ความได้เปรียบทั้งหมดมาจากกลุ่มคำค้นที่คลังไม่มีสินค้าตรงเท่านั้น ส่วนกลุ่มที่คลังมีสินค้าตรงนั้นทุกสถาปัตยกรรมได้คะแนนเท่ากัน นอกจากนี้ยังพบว่าชั้นจัดอันดับเป็นคอขวดที่แท้จริงของระบบ เนื่องจากทำให้สินค้าที่ตรงสูญหายไปร้อยละ 45 ถึง 48 ซึ่งมากกว่าทุกชั้นอื่นรวมกันประมาณ 3.5 เท่า และการเพิ่มสัญญาณนับคำที่ซ้อนกับคำค้นเข้าไปในชั้นนี้ให้ผลดีกว่าการเพิ่มความเป็นเอเจนต์ประมาณ 2 เท่า

ผลที่ได้นำไปใช้เป็นแนวทางเลือกสถาปัตยกรรมสำหรับระบบค้นหาสินค้าได้ คือความเป็นเอเจนต์คุ้มค่าเมื่อคำค้นมีโอกาสหาของไม่เจอ ไม่ใช่ทุกกรณี และโครงงานนี้ได้เผยแพร่ชุดทดสอบพร้อมเฉลยสู่สาธารณะเพื่อให้ผู้อื่นนำไปวัดผลซ้ำได้
